% =============================================================================
% Section 3.4: Background on Video Streaming Technology
% =============================================================================

\section{Background on Video Streaming Technology}
\label{sec:video-streaming}

\subsection{HTTP Live Streaming (HLS)}

HTTP Live Streaming (HLS), standardised as RFC 8216~\cite{pantos2022hls}, is the dominant adaptive-bitrate streaming protocol used by major video platforms including YouTube, Twitch, and Netflix. HLS works by segmenting a source video into short (typically 2--10 second) chunks encoded at multiple bitrates and resolutions. A \textbf{manifest file} (the \texttt{.m3u8} playlist) enumerates the available renditions and their segment URLs. The client---a web browser or mobile app---dynamically selects the highest bitrate segment that its current network conditions can sustain, switching to a lower bitrate when bandwidth drops and back to a higher bitrate when bandwidth recovers. This adaptive behaviour eliminates buffering interruptions while maximising playback quality.

The HLS specification defines a hierarchical playlist structure. A \textit{master playlist} lists all variant streams (e.g., 360p, 720p, 1080p) and auxiliary media such as subtitles and alternate audio tracks. Each variant stream has its own \textit{media playlist} containing the sequence of segment URIs and their durations. The protocol supports encryption via AES-128 or SAMPLE-AES, allowing content access control through key delivery authorisation.

\subsection{Mux Video Platform}

Mux~\cite{mux2024} is a video API platform that abstracts the complexity of video ingestion, transcoding, storage, and delivery. NexaLearn integrates with Mux through the \textbf{Media} bounded context, which manages the full lifecycle of video assets. Mux provides:

\begin{itemize}
  \item \textbf{Direct Uploads.} The client uploads a video file directly to Mux's storage via a \textit{direct upload endpoint}, obtaining a \texttt{upload\_id} and a playback ID. The upload is initiated by NexaLearn's backend, which requests a signed upload URL from Mux and returns it to the client; the client then PUTs the file directly to Mux, avoiding double-hop through the NexaLearn server.
  \item \textbf{Transcoding Pipeline.} Upon receiving a complete upload, Mux transcodes the source video into multiple HLS renditions (typically 360p, 480p, 720p, and 1080p) with configurable video codec (H.264), audio codec (AAC), and segment duration.
  \item \textbf{Webhook Notifications.} Mux sends webhooks to NexaLearn for lifecycle events: \texttt{video.upload.asset\_created}, \texttt{video.asset.ready}, \texttt{video.asset.errored}, \texttt{video.asset.deleted}. These webhooks drive the video asset state machine.
  \item \textbf{Playback IDs and Signed URLs.} Each transcoded asset receives a playback ID. For gated content, Mux supports \textit{signed playback URLs} using JSON Web Tokens (JWT)~\cite{jones2015jwt}. The NexaLearn backend generates a JWT containing the playback ID, an expiration time, and optional access restrictions (e.g., referrer whitelist, IP allowlist), signed with a Mux-specific signing key. The client embeds this signed JWT in the playback URL; Mux's delivery infrastructure verifies the signature and expiration before serving segments.
\end{itemize}

\subsection{Video Asset Lifecycle in NexaLearn}

The Media bounded context implements a state machine for each \texttt{VideoAsset} entity. The states are:

\begin{description}
  \item[\texttt{AWAITING\_UPLOAD}.] Initial state. A \texttt{Video} or \texttt{Lesson} entity has been created with an association to a new video, but the instructor has not yet uploaded the file. No Mux resource exists.
  \item[\texttt{UPLOADED}.] The instructor has initiated an upload and Mux has confirmed receipt of the source file. The state transitions from AWAITING\_UPLOAD to UPLOADED.
  \item[\texttt{TRANSCODING}.] Mux sends \texttt{video.upload.asset\_created}, indicating that the source file is being transcoded into HLS renditions. The NexaLearn backend creates a \texttt{GenerationJob} entity for the AI pipeline (Section~\ref{sec:ai-elearning}) in QUEUED state, awaiting the completed transcript.
  \item[\texttt{READY}.] Mux sends \texttt{video.asset.ready}, indicating that all HLS renditions are available. NexaLearn updates the \texttt{VideoAsset} with the playback ID and signed URL expiration policy. The AI pipeline receives a \texttt{VideoReady} event via the outbox, triggering the transcript-generation workflow.
  \item[\texttt{FAILED}.] Transcoding failed. NexaLearn logs the Mux error code and message, notifies the instructor, and makes the upload available for retry.
\end{description}

\subsection{MP4 Renditions for the AI Pipeline}

While HLS is the delivery format for learner playback, the Whisper transcription model (Section~\ref{sec:ai-elearning}) requires a contiguous audio or video file, not segmented HLS chunks. Mux provides a static MP4 rendition for each asset, accessible through a separate signed URL with an extended expiration window---sufficient for the AI pipeline to download the file, extract the audio track using FFmpeg, and feed it to Whisper. This dual-delivery strategy---HLS for learners, MP4 for automated processing---is transparent to the instructor, who uploads a single source file.

The interplay between video lifecycle and the AI pipeline exemplifies NexaLearn's event-driven design: a single external event (\texttt{video.asset.ready}) propagates through the outbox to trigger transcript generation, quiz creation, instructor notification, and search index updates, all without explicit orchestration code in the Media context. An asset may also transition directly to FAILED from UPLOADED or TRANSCODING if an error occurs at any stage.
